\documentclass{amsart}
% For dealing with references we use the comment environment
\usepackage{verbatim}
\newenvironment{reference}{\comment}{\endcomment}
%\newenvironment{reference}{}{}
\newenvironment{slogan}{\comment}{\endcomment}
\newenvironment{history}{\comment}{\endcomment}

% For commutative diagrams we use Xy-pic
\usepackage[all]{xy}

% We use 2cell for 2-commutative diagrams.
\xyoption{2cell}
\UseAllTwocells

% We use multicol for the list of chapters between chapters
\usepackage{multicol}

% This is generally recommended for better output
\usepackage{lmodern}
\usepackage[T1]{fontenc}

% For cross-file-references

% Package for hypertext links:
\usepackage{hyperref}

% For any local file, say "hello.tex" you want to link to please

% Theorem environments.
%
\theoremstyle{plain}
\newtheorem{theorem}[subsection]{Theorem}
\newtheorem{proposition}[subsection]{Proposition}
\newtheorem{lemma}[subsection]{Lemma}

\theoremstyle{definition}
\newtheorem{definition}[subsection]{Definition}
\newtheorem{example}[subsection]{Example}
\newtheorem{exercise}[subsection]{Exercise}
\newtheorem{situation}[subsection]{Situation}

\theoremstyle{remark}
\newtheorem{remark}[subsection]{Remark}
\newtheorem{remarks}[subsection]{Remarks}

\numberwithin{equation}{subsection}

% Macros
%
\def\lim{\mathop{\mathrm{lim}}\nolimits}
\def\colim{\mathop{\mathrm{colim}}\nolimits}
\def\Spec{\mathop{\mathrm{Spec}}}
\def\Hom{\mathop{\mathrm{Hom}}\nolimits}
\def\Ext{\mathop{\mathrm{Ext}}\nolimits}
\def\SheafHom{\mathop{\mathcal{H}\!\mathit{om}}\nolimits}
\def\SheafExt{\mathop{\mathcal{E}\!\mathit{xt}}\nolimits}
\def\Sch{\mathit{Sch}}
\def\Mor{\mathop{\mathrm{Mor}}\nolimits}
\def\Ob{\mathop{\mathrm{Ob}}\nolimits}
\def\Sh{\mathop{\mathit{Sh}}\nolimits}
\def\NL{\mathop{N\!L}\nolimits}
\def\CH{\mathop{\mathrm{CH}}\nolimits}
\def\proetale{{pro\text{-}\acute{e}tale}}
\def\etale{{\acute{e}tale}}
\def\QCoh{\mathit{QCoh}}
\def\Ker{\mathop{\mathrm{Ker}}}
\def\Im{\mathop{\mathrm{Im}}}
\def\Coker{\mathop{\mathrm{Coker}}}
\def\Coim{\mathop{\mathrm{Coim}}}

% Boxtimes
%
\DeclareMathSymbol{\boxtimes}{\mathbin}{AMSa}{"02}

%
% Macros for moduli stacks/spaces
%
\def\QCohstack{\mathcal{QC}\!\mathit{oh}}
\def\Cohstack{\mathcal{C}\!\mathit{oh}}
\def\Spacesstack{\mathcal{S}\!\mathit{paces}}
\def\Quotfunctor{\mathrm{Quot}}
\def\Hilbfunctor{\mathrm{Hilb}}
\def\Curvesstack{\mathcal{C}\!\mathit{urves}}
\def\Polarizedstack{\mathcal{P}\!\mathit{olarized}}
\def\Complexesstack{\mathcal{C}\!\mathit{omplexes}}
% \Pic is the operator that assigns to X its picard group, usage \Pic(X)
% \Picardstack_{X/B} denotes the Picard stack of X over B
% \Picardfunctor_{X/B} denotes the Picard functor of X over B
\def\Pic{\mathop{\mathrm{Pic}}\nolimits}
\def\Picardstack{\mathcal{P}\!\mathit{ic}}
\def\Picardfunctor{\mathrm{Pic}}
\def\Deformationcategory{\mathcal{D}\!\mathit{ef}}

\setlength{\emergencystretch}{3em}
\begin{document}
\title[Stacks Project, Tag 0FA9]{452 --- Top Chern class equals the zero-cycle class for sections with expected dimension and generic regularity}
\maketitle
\noindent Copyright (C) 2005--2025 Johan de Jong. Stacks Project authors. Source: \href{https://stacks.math.columbia.edu/tag/0FA9}{Tag 0FA9}.

\noindent Editorial difficulty note (not author text). Difficulty H2: The proof must preserve the zero-locus multiplicities while reducing the vector-bundle rank. Pullback to the projective bundle produces a line-bundle quotient section and a lower-rank section on its zero locus; regular-sequence control and fibre-dimension estimates are needed to justify the induction..

\noindent Editorial provenance note (not author text). History: excerpt from revision \texttt{a0444\allowbreak 6e57e\allowbreak c1fbc\allowbreak 252a8\allowbreak 71afc\allowbreak ec775\allowbreak 2fb28\allowbreak 07b14}, chow.tex, lines 8749--8836. Reference presentation was adapted; original-author context and core support blocks were retained or added. The mathematical wording is unchanged. Licensed under GNU FDL 1.2 or later; no invariant sections or cover texts. See ../licenses/COPYING. Context blocks reproduce author definitions and supporting results; they are not additional counted problems.

\section{Chern classes and sections}
\label{section-top-chern-class}

\noindent
A brief section whose main result is that we may compute the
top Chern class of a finite locally free module using the
vanishing locus of a ``regular section.

\medskip\noindent
Let $(S, \delta)$ be as in Situation \ref{situation-setup}. Let $X$ be a scheme
locally of finite type over $S$. Let $\mathcal{E}$ be a finite locally
free $\mathcal{O}_X$-module. Let $f : X' \to X$ be locally of finite type. Let
$$
s \in \Gamma(X', f^*\mathcal{E})
$$
be a global section of the pullback of $\mathcal{E}$ to $X'$. Let
$Z(s) \subset X'$ be the zero scheme of $s$. More precisely, we define
$Z(s)$ to be the closed subscheme whose quasi-coherent sheaf
of ideals is the image of the map
$s : f^*\mathcal{E}^\vee \to \mathcal{O}_{X'}$.



\begin{situation}
\label{situation-setup}
Here $S$ is a locally Noetherian, and universally catenary scheme.
Moreover, we assume $S$ is endowed with a dimension function
$\delta : S \longrightarrow \mathbf{Z}$.
\end{situation}

\begin{lemma}
\label{lemma-top-chern-class}
In the situation described just above assume $\dim_\delta(X') = n$,
that $f^*\mathcal{E}$ has constant rank $r$, that
$\dim_\delta(Z(s)) \leq n - r$, and that for every generic point
$\xi \in Z(s)$ with $\delta(\xi) = n - r$ the ideal of $Z(s)$
in $\mathcal{O}_{X', \xi}$ is generated by a regular sequence
of length $r$. Then
$$
c_r(\mathcal{E}) \cap [X']_n = [Z(s)]_{n - r}
$$
in $\CH_*(X')$.
\end{lemma}

\begin{proof}
Since $c_r(\mathcal{E})$ is a bivariant class
(Lemma \href{https://stacks.math.columbia.edu/tag/0B7H}{lemma cap cp bivariant (Tag 0B7H)})
we may assume $X = X'$ and we have to show that
$c_r(\mathcal{E}) \cap [X]_n = [Z(s)]_{n - r}$ in $\CH_{n - r}(X)$.
We will prove the lemma by induction on $r \geq 0$. (The case
$r = 0$ is trivial.) The case $r = 1$
is handled by Lemma \href{https://stacks.math.columbia.edu/tag/02SQ}{lemma geometric cap (Tag 02SQ)}. Assume $r > 1$.

\medskip\noindent
Let $\pi : P \to X$ be the projective space bundle associated to
$\mathcal{E}$ and consider the short exact sequence
$$
0 \to \mathcal{E}' \to \pi^*\mathcal{E} \to \mathcal{O}_P(1) \to 0
$$
By the projective space bundle formula
(Lemma \href{https://stacks.math.columbia.edu/tag/02TX}{lemma chow ring projective bundle (Tag 02TX)})
it suffices to prove the equality after pulling back by $\pi$.
Observe that $\pi^{-1}Z(s) = Z(\pi^*s)$ has $\delta$-dimension
$\leq n - 1$ and that the assumption on regular sequences at
generic points of $\delta$-dimension $n - 1$ holds by
flat pullback, see
Algebra, Lemma \href{https://stacks.math.columbia.edu/tag/00LM}{lemma flat increases depth (Tag 00LM)}.
Let $t \in \Gamma(P, \mathcal{O}_P(1))$ be the image of $\pi^*s$.
We claim
$$
[Z(t)]_{n + r - 2} = c_1(\mathcal{O}_P(1)) \cap [P]_{n + r - 1}
$$
Assuming the claim we finish the proof as follows.
The restriction $\pi^*s|_{Z(t)}$ maps to zero in
$\mathcal{O}_P(1)|_{Z(t)}$ hence comes from a unique
element $s' \in \Gamma(Z(t), \mathcal{E}'|_{Z(t)})$.
Note that $Z(s') = Z(\pi^*s)$ as closed subschemes of $P$.
If $\xi \in Z(s')$ is a generic point with $\delta(\xi) = n - 1$,
then the ideal of $Z(s')$ in $\mathcal{O}_{Z(t), \xi}$
can be generated by a regular sequence of length $r - 1$: it is generated by
$r - 1$ elements which are the images of $r - 1$ elements in
$\mathcal{O}_{P, \xi}$ which together with a generator of the
ideal of $Z(t)$ in $\mathcal{O}_{P, \xi}$ form a regular sequence
of length $r$ in $\mathcal{O}_{P, \xi}$. Hence we can apply the
induction hypothesis to $s'$ on $Z(t)$ to get
$c_{r - 1}(\mathcal{E}') \cap [Z(t)]_{n + r - 2} = [Z(s')]_{n - 1}$.
Combining all of the above we obtain
\begin{align*}
c_r(\pi^*\mathcal{E}) \cap [P]_{n + r - 1}
& =
c_{r - 1}(\mathcal{E}') \cap c_1(\mathcal{O}_P(1)) \cap [P]_{n + r - 1} \\
& =
c_{r - 1}(\mathcal{E}') \cap [Z(t)]_{n + r - 2} \\
& =
[Z(s')]_{n - 1} \\
& = [Z(\pi^*s)]_{n - 1}
\end{align*}
which is what we had to show.

\medskip\noindent
Proof of the claim. This will follow from an application of
the already used Lemma \href{https://stacks.math.columbia.edu/tag/02SQ}{lemma geometric cap (Tag 02SQ)}.
We have $\pi^{-1}(Z(s)) = Z(\pi^*s) \subset Z(t)$.
On the other hand, for $x \in X$ if $P_x \subset Z(t)$, then
$t|_{P_x} = 0$ which implies that $s$ is zero in the fibre
$\mathcal{E} \otimes \kappa(x)$, which implies $x \in Z(s)$.
It follows that $\dim_\delta(Z(t)) \leq n + (r - 1) - 1$.
Finally, let $\xi \in Z(t)$ be a generic point with
$\delta(\xi) = n + r - 2$. If $\xi$ is not the generic point
of the fibre of $P \to X$ it is immediate that
a local equation of $Z(t)$ is a nonzerodivisor in $\mathcal{O}_{P, \xi}$
(because we can check this on the fibre by
Algebra, Lemma \href{https://stacks.math.columbia.edu/tag/00MF}{lemma grothendieck (Tag 00MF)}).
If $\xi$ is the generic point of a fibre, then $x = \pi(\xi) \in Z(s)$
and $\delta(x) = n + r - 2 - (r - 1) = n - 1$. This is a contradiction
with $\dim_\delta(Z(s)) \leq n - r$ because $r > 1$
so this case doesn't happen.
\end{proof}
\end{document}
